\documentclass[11pt]{article}
\usepackage{amsmath,amssymb}
\usepackage[margin=1in]{geometry}

\newcommand{\rhoN}{\rho_{1/p}}
\newcommand{\deltaN}{\delta_N}

\title{Step 195: Beurling--Nyman Arithmetic Gram Refinement}
\date{}

\begin{document}
\maketitle

\section*{Verdict}

\[
\boxed{\texttt{V\_BN\_arithmetic\_refined\_RH\_consistent}}.
\]

The harmonic Beurling--Nyman chain \(A_N=\{1/k:1\le k\le N\}\) was extended to \(N=1000\).  The refined data remain consistent with the slow logarithmic decay expected in the Baez--Duarte/Nyman--Beurling framework.

\section*{Arithmetic Gram Formula}

For \(a=1/p\), after the substitution \(x=1/t\),
\[
\rho_{1/p}(1/x)
=
\frac{\lfloor x\rfloor}{p}-\left\lfloor \frac{x}{p}\right\rfloor.
\]
On the unit interval \(x\in(n,n+1)\), this equals
\[
\frac{n\bmod p}{p}.
\]
Therefore
\[
G_{p,q}
=
\langle \rho_{1/p},\rho_{1/q}\rangle
=
\sum_{n\ge1}
\frac{(n\bmod p)(n\bmod q)}{pq\,n(n+1)}.
\]
Using periodicity with \(L=\operatorname{lcm}(p,q)\), this is equivalently
\[
G_{p,q}
=
\frac1L\sum_{r=1}^L
\left\{\frac rp\right\}\left\{\frac rq\right\}
\left[
\psi\!\left(\frac{r+1}{L}\right)
-
\psi\!\left(\frac rL\right)
\right],
\]
where \(\psi\) is the digamma function.  The implementation spot-checks this periodized formula at 80 decimal digits for \((p,q)=(2,3),(3,5),(7,11)\).

For the dense \(N=1000\) computation, the script uses the arithmetic series in vectorized form with \(X_{\max}=200000\).  This is no longer a discontinuous integral quadrature.  The remaining arithmetic tail is bounded entrywise by \(5\cdot 10^{-6}\).

\section*{Precision}

\begin{itemize}
\item \(\texttt{mpmath}\) precision: 80 decimal digits.
\item Source vector \(b_p=(1/p)\log p\): computed using \(\texttt{mpmath}\).
\item Dense pseudoinverse: NumPy float64 SVD, \(\texttt{rcond}=10^{-12}\).
\item Tail bound per Gram entry: \(5\cdot10^{-6}\).
\end{itemize}

\section*{Computed Harmonic Defects}

\begin{center}
\begin{tabular}{rcc}
\(N\) & \(\delta_N^2\) & \(\delta_N^2\log N\)\\
\hline
5 & 0.0363139368590198 & 0.0584450267306446\\
10 & 0.0238470859060593 & 0.0549099445186406\\
20 & 0.0165323317860087 & 0.0495264398884486\\
40 & 0.0127296188449689 & 0.0469580294159074\\
80 & 0.0109920955378177 & 0.0481676554175969\\
200 & 0.00892242600836957 & 0.0472738446718843\\
500 & 0.00741541353971342 & 0.0460838890370526\\
1000 & 0.00654423096394308 & 0.0452059459880561\\
\end{tabular}
\end{center}

\section*{Decay Fit}

For the tail values \(N=80,200,500,1000\),
\[
\operatorname{mean}(\delta_N^2\log N)
=
0.0466828337786475.
\]
A least-squares fit over \(N\ge20\) to
\[
\delta_N^2\sim \frac{C}{(\log N)^\alpha}
\]
gives
\[
C=0.0541389792662914,
\qquad
\alpha=1.0886679.
\]
The comparison-only power fit \(\delta_N^2\sim C N^{-\beta}\) gives
\[
\beta=0.22951419,
\]
but the logarithmic model is the Beurling--Nyman/Baez--Duarte relevant one.

\section*{Comparison to Baez--Duarte}

Baez--Duarte's strengthening makes the \(1/k\) chain a legitimate countable carrier for the criterion, and the numerical literature around this criterion predicts very slow logarithmic-scale decay rather than fast algebraic convergence.  The \(N\le1000\) data are consistent with that slow-log regime.  They do not establish the limiting statement.

\section*{Nonclaim}

This is a finite numerical refinement.  It does not establish RH and does not close the Beurling--Nyman residual.

\end{document}

